\documentclass[11pt]{article}
\usepackage[margin=1in]{geometry}
\usepackage{amsmath,amssymb,amsthm,mathtools}
\usepackage[colorlinks=true,linkcolor=blue,citecolor=blue,urlcolor=blue]{hyperref}
\usepackage{enumitem}
\newtheorem{theorem}{Theorem}[section]
\newtheorem{proposition}[theorem]{Proposition}
\newtheorem{lemma}[theorem]{Lemma}
\newtheorem{corollary}[theorem]{Corollary}
\theoremstyle{definition}
\newtheorem{definition}[theorem]{Definition}
\theoremstyle{remark}
\newtheorem{remark}[theorem]{Remark}
\newcommand{\CP}{\mathbb{CP}}
\newcommand{\C}{\mathbb C}
\newcommand{\Ric}{\operatorname{Ric}}
\newcommand{\Riem}{\operatorname{Riem}}
\newcommand{\End}{\operatorname{End}}
\newcommand{\Sym}{\operatorname{Sym}}
\newcommand{\id}{\operatorname{id}}
\newcommand{\dd}{\mathrm d}
\newcommand{\rank}{\operatorname{rank}}

\title{Full-connection descent and chiral curvature nullity\\
\large Proposed section with derivation and hypothesis audit}
\author{Working research notes}
\date{21 November 2026}

\begin{document}
\maketitle

\noindent\textbf{Status.}
The obstruction below is derived in these notes from curvature horizontality
and two-spinor linear algebra. It is not attributed to a published theorem,
and no claim of bibliographic novelty is made. The cited Eguchi--Hanson
calculation supplies an example to which the derived obstruction applies.
The notes distinguish this full-connection statement from the weaker
transfer of alpha-partial connections or auxiliary Ward bundles.

\section{Conventions and the meaning of full descent}

All manifolds, bundles, connections, and maps are holomorphic and local.
Let $(M,g)$ be a four-complex-dimensional spin metric manifold, with
\[
 TM\simeq S\otimes S',\qquad \rank S=\rank S'=2.
\]
The spin bundles carry symplectic forms $\epsilon_{AB}$ and
$\epsilon_{A'B'}$. An alpha plane is a totally null plane of the form
$S_x\otimes\langle\pi^{A'}\rangle$; its unprimed factor varies and its
primed factor is fixed. Let $\gamma,\gamma'$ be the Levi--Civita spin
connections and write
\[
 F^S=\dd\gamma+\gamma\wedge\gamma,
 \qquad F^{S'}=\dd\gamma'+\gamma'\wedge\gamma'.
\]
Here $F^S$ denotes spin curvature, $\Riem$ denotes the curvature of the
Levi--Civita connection on $TM$, and $\Ric$ denotes its Ricci contraction.
These three objects are not denoted by the same symbol.

We use the chirality convention
\begin{equation}\label{obs:eq:vacuum-asd}
 \Ric(g)=0,\qquad \widetilde\Psi_{A'B'C'D'}=0.
\end{equation}
Then $F^{S'}=0$, while the potentially nonzero curvature is
\begin{equation}\label{obs:eq:spin-curvature}
 (F^S)_{AA'BB'CD}=\epsilon_{A'B'}\Psi_{ABCD},
 \qquad \Psi_{ABCD}=\Psi_{(ABCD)}.
\end{equation}
The last endomorphism index has been lowered with the spin symplectic form.
An alternative curvature or lowering convention can multiply the right
side by a nonzero constant; none of the nullity statements changes.
As endomorphism-valued two-forms,
\begin{equation}\label{obs:eq:riemann-spin}
 \Riem=F^S\otimes\id_{S'}+\id_S\otimes F^{S'}
       =F^S\otimes\id_{S'}.
\end{equation}
In particular, under \eqref{obs:eq:vacuum-asd}, vanishing $F^S$ is equivalent
to vanishing $\Riem$.

Let $\varpi:\mathbb P(S')\to M$ be the curved correspondence projection,
and let $D_g$ be the rank-two alpha distribution. The projection
$\dd\varpi$ identifies $D_g$ at $(x,[\pi])$ with the alpha plane
$S_x\otimes\langle\pi\rangle$. On a flat four-dimensional neighborhood
$(U,\eta)$, use the flat correspondence space $U\times\CP^1$, projection
$p$, and its standard distribution $D_\eta$.

Suppose
\begin{equation}\label{obs:eq:inverse-chart}
 \Phi:W\longrightarrow\mathbb P(S'),
 \qquad W\subset U\times\CP^1,
\end{equation}
is a local biholomorphism satisfying
$\dd\Phi(D_\eta)=D_g$. Put $r=\varpi\circ\Phi$. No preservation of either
spacetime projection is assumed. After shrinking around a chosen point,
$W$ can be taken to be a product $U_0\times D$, where $D$ is a direction
disc with coordinate $\zeta$.

\begin{definition}\label{obs:def:full-descent}
The full spin connection descends through $\Phi$ if there are a holomorphic
bundle $E_0\to U_0$ with connection $\nabla^A$ and a bundle isomorphism over
$W$ that identifies the entire connection $r^*\nabla^S$ with
$p^*\nabla^A$. Equivalently, in local frames, $r^*\gamma$ is gauge-equivalent
to a one-form $A_i(y)\dd y^i$ with no direction component and no direction
dependence. The gauge transformation may depend on $(y,\zeta)$.
\end{definition}

This is equality of full connections in all five correspondence-space
directions. Equality only after restriction to $D_\eta$ is a different and
weaker requirement.

\section{The exact local gauge criterion}

\begin{lemma}[Local descent up to gauge]\label{obs:lem:basicness}
Let $\nabla^B$ be a holomorphic connection on a vector bundle over a product
$U_0\times D$. After shrinking around any chosen point, it is
gauge-equivalent to the pullback of a connection on $U_0$ if and only if
\begin{equation}\label{obs:eq:horizontal-curvature}
 \iota_V F_B=0\qquad\text{for every }V\in\ker\dd p.
\end{equation}
The condition includes all mixed curvature components, not just the
restriction of the curvature to two vertical vectors.
\end{lemma}

\begin{proof}
Necessity follows because the curvature of a pullback is $p^*F_A$ and
curvature changes by conjugation under a gauge transformation.

For sufficiency, choose a local frame and write
\[
 B=B_\zeta(y,\zeta)\dd\zeta+\sum_i B_i(y,\zeta)\dd y^i.
\]
Solve the holomorphic matrix ordinary differential equation with parameters
\begin{equation}\label{obs:eq:vertical-ode}
 \partial_\zeta h=-B_\zeta h,
 \qquad h(y,\zeta_0)=I.
\end{equation}
After shrinking, its solution is holomorphic and invertible. With the gauge
convention
\[
 B^h=h^{-1}Bh+h^{-1}\dd h,
\]
equation \eqref{obs:eq:vertical-ode} gives $B^h_\zeta=0$.
Condition \eqref{obs:eq:horizontal-curvature} is gauge-invariant, hence
\[
 0=(F_{B^h})_{\zeta i}=\partial_\zeta B^h_i.
\]
Thus $B^h_i=A_i(y)$, as required. If $B$ is traceless, the initial condition
also gives $\det h=1$, so the proof can use $SL(2,\C)$ gauges for the spin
connection rather than requiring the larger group $GL(2,\C)$.
\end{proof}

\begin{remark}
The lemma is local on a direction disc. It is not a statement that a chosen
gauge extends holomorphically over a complete $\CP^1$, nor that a single
inverse correspondence chart exists over that fiber. Those global
questions require additional hypotheses. Also, in complex dimension one
the curvature on pairs of vertical vectors vanishes automatically; that
vacuous condition is much weaker than \eqref{obs:eq:horizontal-curvature}.
\end{remark}

\begin{proposition}[Curvature-nullity criterion for a fixed chart]
\label{obs:prop:fixed-chart}
For the map \eqref{obs:eq:inverse-chart}, define
\begin{equation}\label{obs:eq:nullity}
 \mathcal K_x=\{v\in T_xM:F^S_x(v,w)=0\text{ for all }w\in T_xM\}.
\end{equation}
Locally, full descent through $\Phi$ is equivalent to
\begin{equation}\label{obs:eq:vertical-nullity}
 \dd r(\ker\dd p)\subset\mathcal K.
\end{equation}
No constant-rank hypothesis on $\mathcal K$ is needed for this equivalence.
\end{proposition}

\begin{proof}
Because $\Phi$ is a local biholomorphism and $\varpi$ is a submersion,
$\dd r=\dd\varpi\circ\dd\Phi$ is surjective. The curvature of the full
pullback is $r^*F^S$, so for a target-vertical vector $V$ and any $Y\in TW$,
\[
 (r^*F^S)(V,Y)=F^S(\dd r(V),\dd r(Y)).
\]
Surjectivity of $\dd r$ makes its vanishing for every $Y$ equivalent to
$\dd r(V)\in\mathcal K$. Apply Lemma~\ref{obs:lem:basicness}.
\end{proof}

\section{Chiral curvature nullity and the local obstruction}

\begin{lemma}[Pointwise spinor nullity]\label{obs:lem:spinor-nullity}
Under \eqref{obs:eq:vacuum-asd}, let
\[
 T_\Psi:S_x\longrightarrow\Sym^3S_x^*,
 \qquad (T_\Psi\lambda)_{BCD}=\lambda^A\Psi_{ABCD}.
\]
Then
\begin{equation}\label{obs:eq:kernel-tensor}
 \mathcal K_x=(\ker T_\Psi)\otimes S'_x.
\end{equation}
Consequently:
\begin{enumerate}[label=\textup{(\roman*)}]
\item if $\Psi_x=0$, then $\dim_\C\mathcal K_x=4$;
\item if $\Psi_x=\psi\,\kappa_A\kappa_B\kappa_C\kappa_D\ne0$
      (Petrov type N), then $\dim_\C\mathcal K_x=2$;
\item for every other nonzero algebraic type, $\mathcal K_x=0$.
\end{enumerate}
In case \textup{(ii)}, $\mathcal K_x=(\ker\kappa)\otimes S'_x$ is totally
null. With our convention for alpha planes, it belongs to the opposite
ruling, often called the beta ruling.
\end{lemma}

\begin{proof}
Contracting \eqref{obs:eq:spin-curvature} in the first tangent argument
shows that $v\in\mathcal K_x$ exactly when
\[
 v^{AA'}\epsilon_{A'B'}\Psi_{ABCD}=0
 \qquad\text{for every }B',B,C,D.
\]
Invertibility of $\epsilon_{A'B'}$ says precisely that each column
$v^{AA'}$, viewed as an unprimed spinor, lies in $\ker T_\Psi$.
This proves \eqref{obs:eq:kernel-tensor}.

If $0\ne\lambda\in\ker T_\Psi$, choose a basis of $S_x$ with
$\lambda=e_1$. Then $\Psi_{1BCD}=0$. Complete symmetry implies that every
component with any index equal to $1$ vanishes. The only possible nonzero
component is $\Psi_{0000}$, so $\Psi=\psi\,(e^0)^4$. Conversely, a nonzero
fourth power has the one-dimensional kernel $\ker\kappa$. If
$\dim\ker T_\Psi=2$, all components vanish. Tensoring with $S'_x$ gives
the stated dimensions. The plane $(\ker\kappa)\otimes S'_x$ is totally
null since its unprimed factors are proportional and $\epsilon$ is
alternating.
\end{proof}

\begin{lemma}[An open family of alpha planes determines the conformal class]
\label{obs:lem:alpha-conformal}
Let $f$ be a local biholomorphism from a flat four-manifold $(U_0,\eta)$
to $(M,g)$. Suppose that, for every $y$, $\dd f_y$ maps the flat alpha
planes in a nonempty open set of projective directions to totally null
planes for $g$. Then $f^*g=c\eta$ for a nowhere-zero holomorphic function
$c$.
\end{lemma}

\begin{proof}
At fixed $y$, write a flat null vector as $v^{AA'}=\lambda^A\pi^{A'}$.
The quadratic form $q_y(v)=(f^*g)_y(v,v)$ vanishes for every $\lambda$ and
for $[\pi]$ in the specified open set. As a polynomial, $q_y(\lambda\pi)$
has degree two in each spinor and therefore vanishes for all $\lambda,\pi$.
Thus $q_y$ vanishes on the entire flat null quadric. The determinant
polynomial defining that quadric is irreducible over $\C$, so the two
quadratic forms are proportional. Nondegeneracy gives $c(y)\ne0$.
Comparing any fixed nonzero component of $\eta$ makes $c$ holomorphic.
\end{proof}

\begin{theorem}[Local algebraic obstruction to full spin-connection descent]
\label{obs:thm:petrov-obstruction}
Assume \eqref{obs:eq:vacuum-asd} and let $\Phi$ satisfy
\eqref{obs:eq:inverse-chart}, including alpha-distribution preservation.
If the full connection $r^*\nabla^S$ descends in the sense of
Definition~\ref{obs:def:full-descent}, then at every point of $r(W)$ the
Weyl spinor $\Psi$ is of Petrov type N or is zero (type O).

The conclusion allows arbitrary holomorphic direction-dependent gauges,
requires no complete projective direction fiber, and does not require the
descended connection to be stipulated ASD in advance.
\end{theorem}

\begin{proof}
Suppose, to the contrary, that $\Psi$ has a nonzero type other than N at
$x_0=r(y_0,\zeta_0)$. Lemma~\ref{obs:lem:spinor-nullity} says that
$T_\Psi$ is injective at $x_0$. Injectivity is an open condition on a
holomorphic matrix (some two-by-two minor is nonzero), so
$\mathcal K=0$ on an open neighborhood $M_0$ of $x_0$.

Shrink to a product neighborhood $U_0\times D$ of $(y_0,\zeta_0)$ whose
image under $r$ lies in $M_0$. Proposition~\ref{obs:prop:fixed-chart} gives
$\partial_\zeta r=0$. Hence
\[
 r(y,\zeta)=f(y).
\]
Since $r$ is a submersion, $\dd f$ has rank four; shrinking again makes
$f$ a local biholomorphism.

For each $(y,\zeta)$, the equality
$\dd\Phi(D_\eta)=D_g$ and the identity $\varpi\circ\Phi=f\circ p$
imply that $\dd f_y$ maps the flat alpha plane labelled by $\zeta$ to the
source alpha plane labelled by $\Phi(y,\zeta)$. The disc $D$ supplies a
nonempty open family of flat directions. Lemma~\ref{obs:lem:alpha-conformal}
gives $f^*g=c\eta$, so $g$ is conformally flat on $f(U_0)$.
Its Weyl tensor therefore vanishes there, in particular at $x_0$,
contradicting the choice of $x_0$. Since no such point exists, the only
possibilities throughout $r(W)$ are N and O.
\end{proof}

\begin{remark}[Scope of the conclusion]
Type N is necessary, not sufficient. The argument does not produce an
alpha-preserving chart with its vertical motion in $\mathcal K$, a
parallel frame of that null plane, or one-dimensional holonomy. Nor does
it assert that Petrov type N cannot change to O on a zero locus. The proof
uses openness only near a point where $T_\Psi$ is injective; it does not
divide by a Weyl component or assume a globally constant Petrov type.
\end{remark}

\begin{remark}[Why the selected chirality matters]
The theorem concerns the potentially nonflat connection $\gamma$ on $S$.
The other connection has $F^{S'}=0$ under our hypotheses and can be made
locally zero. Transporting only that flat chiral connection cannot retain
the gravitational curvature and does not give the stated obstruction.
Exchanging the primed and unprimed conventions exchanges all the chiral
labels but leaves the theorem unchanged.
\end{remark}

\section{The wave model and Eguchi--Hanson}

For the manuscript's wave metric
\[
 g_H=2\dd u\dd v+2\dd z\dd w+H(u,z)\dd u^2,
\]
the nonflat spin curvature is
\[
 F^S=\tfrac12H_{zz}\,\dd z\wedge\dd u\,N,
 \qquad N=\begin{pmatrix}0&0\\1&0\end{pmatrix}.
\]
Where $H_{zz}\ne0$, its common tangent nullity is
$\mathcal K=\operatorname{span}\{\partial_v,\partial_w\}$.
The inverse generalized charts recover $u=-b$ and $z=a$ independently of
direction, so their vertical motion lies exactly in that nullity plane.
The local criterion explains why the full spin-connection pullback can
descend in this example. Where $H_{zz}=0$, the pointwise nullity becomes
larger, and the displayed construction remains regular.

\begin{corollary}[Eguchi--Hanson excludes literal full descent]
\label{obs:cor:eh}
On a nonflat Eguchi--Hanson neighborhood, after choosing chirality as in
\eqref{obs:eq:vacuum-asd}, no alpha-preserving local inverse correspondence
chart can make the full potentially nonflat Levi--Civita spin connection
gauge-equivalent to a connection pulled back from flat spacetime.
\end{corollary}

\begin{proof}
The nonzero Weyl spinor of Eguchi--Hanson is type D. For an explicit check,
Luna et al.~\cite[Section 5.1, equation (5.8)]{LMNO} write its Weyl quartic,
in the opposite chiral notation, as a nonzero multiple of
\[
 \ell_1(\xi)^2\ell_2(\xi)^2,
\]
where, on their regular chart with $\Delta=uv-XY\ne0$,
\[
 \ell_1=-iX(\xi^1-\xi^2)+v(\xi^1+\xi^2),\qquad
 \ell_2=u(\xi^1-\xi^2)+iY(\xi^1+\xi^2).
\]
The determinant of the two coefficient rows is $-2\Delta$, so the two
linear factors are distinct. For their nonzero deformation parameter,
this is type D rather than N or O. Apply
Theorem~\ref{obs:thm:petrov-obstruction}.
\end{proof}

This result does not obstruct the known Eguchi--Hanson Kerr--Schild or Weyl
double copies. It excludes the particular identification of the full
curved spin connection with a basic flat-side connection through an
alpha-preserving correspondence-space local biholomorphism. Transfer of
an auxiliary Ward bundle, or comparison only along the alpha distribution,
does not impose that full identification and remains outside the no-go
statement.

\section{Internal audit and integration notes}

\begin{enumerate}[leftmargin=*,label=\arabic*.]
\item \textbf{Rank is explicit.}
The proof uses local biholomorphicity only to obtain
$\rank\dd r=4$ and, after $\partial_\zeta r=0$, $\rank\dd f=4$.
A lower-rank or singular correspondence map is not covered.

\item \textbf{The relevant curvature is a full two-form.}
The condition is $F(\partial_\zeta,Y)=0$ for every $Y$, not merely
$F|_{D_\eta\wedge D_\eta}=0$. Dropping the mixed components would remove
the obstruction.

\item \textbf{Gauge freedom cannot remove the obstruction.}
It conjugates $F(\partial_\zeta,Y)$, so its vanishing is invariant even
under direction-dependent gauges. The ODE lemma also proves local
sufficiency, not only necessity, for a fixed correspondence map.

\item \textbf{No full-fiber or compactness assumption.}
The theorem needs only a direction disc. It is logically separate from
the compactness obstruction to a full inverse chart over $\CP^1$.

\item \textbf{Reality is unused.}
The proof is complex holomorphic. Petrov type here means the factorization
type of a binary quartic, not a restriction to a Lorentzian real slice.

\item \textbf{Vacuum is essential to the spinor reduction.}
$\Ric=0$ and the vanishing opposite chirality are what reduce spin
curvature to \eqref{obs:eq:spin-curvature}. Nonzero Ricci or cosmological
constant terms require a separate nullity calculation.

\item \textbf{Conformal flatness is sufficient for the contradiction.}
Alpha-plane preservation gives conformal equivalence to the flat metric,
so the Weyl tensor is zero. Vacuum additionally gives $\Riem=0$, but the
contradiction already follows from the nonzero Weyl spinor.

\item \textbf{An open direction family suffices.}
Polynomial continuation on the null quadric replaces any appeal to all
of $\CP^1$. The source direction coordinate need not equal the target
coordinate $\zeta$.

\item \textbf{Zero loci are allowed.}
There is no division by curvature and no assumption that $\mathcal K$
is a bundle globally. The constant-rank step is used only on the open
injectivity locus of $T_\Psi$.

\item \textbf{Pointwise type N is weaker than the wave holonomy class.}
A type-N spinor has a one-dimensional algebraic kernel, but its kernel
line need not be parallel. The theorem does not prove strict Walker
normal form or classify the admissible type-N metrics.

\item \textbf{The descended connection need not be prescribed ASD.}
The impossibility statement allows any ordinary base connection. In the
given setting, alpha-flatness of $F^S$ and alpha preservation would in
fact make a descended $A$ ASD, by the same polynomial-in-direction
argument; this extra fact is not needed for the no-go proof.

\item \textbf{Auxiliary flat Ward bundles are not excluded.}
Line-triviality and a flat alpha-partial connection can reconstruct an
ordinary flat-side ASD connection without identifying its full pullback
with $r^*\gamma$. The resulting connection may therefore survive when
the full-pullback criterion fails.

\item \textbf{Eguchi--Hanson citation and chirality.}
The precise source is Luna--Monteiro--Nicholson--O'Connell,
\emph{Classical and Quantum Gravity} \textbf{36} (2019), 065003,
doi:10.1088/1361-6382/ab03e6, Section 5.1, equation (5.8).
The paragraph following equations (5.24)--(5.25) explicitly identifies
the two double principal spinors. Their nonzero component is written
with dotted indices; our convention exchanges the two chiral labels.
The source was checked against its peer-reviewed repository PDF and
arXiv:1810.08183v3. Its current v3 is dated 1 June 2026; the publication
year remains 2019.

\item \textbf{Status of the derived theorem.}
The proof is supplied in full so it can be checked independently. An
eventual manuscript should call this a proposition or theorem proved
there, not cite it as an established literature result or describe it as
a demonstrated novel theorem without a further literature review.
\end{enumerate}

\begin{thebibliography}{9}
\bibitem{LMNO}
A.~Luna, R.~Monteiro, I.~Nicholson and D.~O'Connell,
\emph{Type D spacetimes and the Weyl double copy},
Class. Quantum Grav. \textbf{36} (2019), 065003.
\href{https://doi.org/10.1088/1361-6382/ab03e6}{doi:10.1088/1361-6382/ab03e6}.
\href{https://arxiv.org/abs/1810.08183}{arXiv:1810.08183}.
\end{thebibliography}

\end{document}
